\documentclass[12pt, a4paper]{article}

% ─── Packages ───
\usepackage[margin=2.5cm]{geometry}
\usepackage{graphicx}
\usepackage{xcolor}
\usepackage{hyperref}
\usepackage{listings}
\usepackage{booktabs}
\usepackage{enumitem}
\usepackage{fancyhdr}
\usepackage{titlesec}
\usepackage{tcolorbox}
\usepackage{tabularx}
\usepackage{parskip}

% ─── Colors ───
\definecolor{orange}{HTML}{F97415}
\definecolor{darkslate}{HTML}{0F172A}
\definecolor{gray}{HTML}{64748B}
\definecolor{codebg}{HTML}{F1F5F9}
\definecolor{green}{HTML}{10B981}

% ─── Hyperlinks ───
\hypersetup{
    colorlinks=true,
    linkcolor=orange,
    urlcolor=orange,
}

% ─── Code Listings ───
\lstset{
    language=Python,
    basicstyle=\ttfamily\small,
    backgroundcolor=\color{codebg},
    frame=single,
    rulecolor=\color{codebg},
    breaklines=true,
    breakatwhitespace=true,
    tabsize=4,
    showstringspaces=false,
    keywordstyle=\color{blue!70!black}\bfseries,
    commentstyle=\color{green!50!black}\itshape,
    stringstyle=\color{orange!80!black},
    xleftmargin=8pt,
    xrightmargin=8pt,
    framexleftmargin=4pt,
    framexrightmargin=4pt,
    aboveskip=10pt,
    belowskip=10pt,
}

% ─── Section Styling ───
\titleformat{\section}
    {\Large\bfseries\color{darkslate}}
    {\thesection.}{0.5em}{}
    [\vspace{-0.5em}{\color{orange}\rule{4cm}{2pt}}]

\titleformat{\subsection}
    {\large\bfseries\color{orange}}
    {\thesubsection}{0.5em}{}

\titleformat{\subsubsection}
    {\normalsize\bfseries\color{darkslate}}
    {\thesubsubsection}{0.5em}{}

% ─── Header / Footer ───
\pagestyle{fancy}
\fancyhf{}
\fancyhead[L]{\small\color{gray}Petpooja AI Copilot}
\fancyhead[R]{\small\color{gray}ML Model Implementation Guide}
\fancyfoot[C]{\small\color{gray}\thepage}
\renewcommand{\headrulewidth}{0.5pt}
\renewcommand{\headrule}{\hbox to\headwidth{\color{orange}\leaders\hrule height \headrulewidth\hfill}}

% ─── Info Box ───
\tcbuselibrary{skins}
\newtcolorbox{infobox}[1][]{
    colback=codebg, colframe=orange, 
    fonttitle=\bfseries, title=#1,
    boxrule=0.5pt, arc=4pt,
    left=8pt, right=8pt, top=6pt, bottom=6pt
}

% ═══════════════════════════════════════════════
\begin{document}

% ─── Title Page ───
\begin{titlepage}
    \centering
    \vspace*{3cm}
    
    {\Huge\bfseries\color{darkslate} Petpooja AI Copilot \par}
    \vspace{0.5cm}
    {\LARGE\color{gray} ML Model Implementation \& Training Guide \par}
    \vspace{1cm}
    {\color{orange}\rule{6cm}{2pt}\par}
    \vspace{1.5cm}
    
    {\large\color{gray}
    A comprehensive technical guide covering four proposed ML models\\[6pt]
    for the Petpooja restaurant intelligence platform:\\[12pt]
    \textbf{Demand Forecasting} $\cdot$ \textbf{NLP Intent Parsing}\\[4pt]
    \textbf{Collaborative Filtering} $\cdot$ \textbf{Anomaly Detection}\\[24pt]
    Each section includes: architecture overview, data requirements,\\
    step-by-step training code, evaluation metrics, and integration strategy.
    \par}
    
    \vfill
    {\color{gray}\today}
\end{titlepage}

% ─── Table of Contents ───
\tableofcontents
\newpage


% ════════════════════════════════════════════════════════════════════
\section{Demand Forecasting (Time-Series Prediction)}
% ════════════════════════════════════════════════════════════════════

\begin{infobox}[Model Overview]
\begin{tabularx}{\textwidth}{@{}lX@{}}
\textbf{Goal}       & Predict how many units of each menu item will sell in the next N days \\
\textbf{Algorithm}  & Facebook Prophet (recommended) / LSTM / ARIMA \\
\textbf{Library}    & \texttt{prophet}, \texttt{pandas}, \texttt{numpy} \\
\textbf{Priority}   & \textcolor{red}{\textbf{HIGH}} --- connects Orders $\to$ Recipes $\to$ Inventory $\to$ Predictions \\
\end{tabularx}
\end{infobox}


\subsection{How It Works}

Time-series forecasting models analyze historical order patterns to predict future demand. The model learns:

\begin{itemize}[nosep]
    \item \textbf{Weekly seasonality} --- weekday vs.\ weekend traffic patterns
    \item \textbf{Daily seasonality} --- lunch rush vs.\ dinner rush vs.\ off-hours
    \item \textbf{Trend} --- overall growth or decline in item popularity
    \item \textbf{External regressors} --- holidays, weather, local events (optional)
\end{itemize}

The output directly feeds into the inventory system: \textit{``You will need 15~kg chicken and 200 burger buns by Friday.''}


\subsection{Data Requirements}

\subsubsection{Minimum Data}
At least \textbf{30 days} of timestamped order history. Accuracy improves significantly with 6+ months of data.

\subsubsection{Data Schema (from \texttt{orders.json})}
\begin{lstlisting}
{
  "order_id": "ORD-20260301-001",
  "timestamp": "2026-03-01T12:34:00",
  "items": [
    { "item_id": "M001", "name": "Classic Veg Burger", "qty": 2 },
    { "item_id": "M019", "name": "French Fries", "qty": 1 }
  ]
}
\end{lstlisting}

\subsubsection{Feature Engineering}
Derive the following features from the raw timestamp:
\begin{itemize}[nosep]
    \item \texttt{day\_of\_week} (0=Mon, 6=Sun) --- captures weekly patterns
    \item \texttt{hour\_of\_day} (0--23) --- captures meal-time peaks
    \item \texttt{is\_weekend} (boolean) --- weekend traffic is typically 30--40\% higher
    \item \texttt{is\_holiday} (boolean) --- from a public holiday calendar
    \item \texttt{rolling\_avg\_7d} --- 7-day rolling average quantity per item
    \item \texttt{weather\_temp} (optional) --- from a weather API; affects beverage sales
\end{itemize}


\subsection{Step-by-Step Training}

\subsubsection{Step 1: Install Dependencies}
\begin{lstlisting}[language=bash]
pip install prophet pandas numpy scikit-learn
\end{lstlisting}

\subsubsection{Step 2: Prepare Training Data}
\begin{lstlisting}
import pandas as pd
import json

# Load historical orders
with open("data/orders.json") as f:
    orders = json.load(f)

# Flatten to daily item counts
rows = []
for order in orders:
    dt = pd.to_datetime(order["timestamp"]).date()
    for item in order["items"]:
        rows.append({
            "ds": dt,
            "item_id": item["item_id"],
            "y": item["qty"]
        })

df = pd.DataFrame(rows)
daily = df.groupby(["ds", "item_id"])["y"].sum().reset_index()
\end{lstlisting}

\subsubsection{Step 3: Train a Prophet Model Per Item}
\begin{lstlisting}
from prophet import Prophet
import pickle

models = {}
for item_id in daily["item_id"].unique():
    item_data = daily[daily["item_id"] == item_id][["ds", "y"]]
    
    model = Prophet(
        daily_seasonality=True,
        weekly_seasonality=True,
        yearly_seasonality=False,   # need 2+ years of data
        changepoint_prior_scale=0.05
    )
    model.fit(item_data)
    models[item_id] = model

# Save all trained models
with open("models/demand_models.pkl", "wb") as f:
    pickle.dump(models, f)
\end{lstlisting}

\subsubsection{Step 4: Generate Predictions}
\begin{lstlisting}
# Predict next 7 days for each item
for item_id, model in models.items():
    future = model.make_future_dataframe(periods=7)
    forecast = model.predict(future)
    
    pred = forecast.tail(7)[["ds", "yhat", "yhat_lower", "yhat_upper"]]
    print(f"{item_id}: next 7 days forecast")
    print(pred.to_string(index=False))
\end{lstlisting}

\subsubsection{Step 5: Connect to Inventory}
\begin{lstlisting}
# For each predicted item demand, calculate ingredient needs
recipes = load_recipes()       # from recipes.json
inventory = load_inventory()   # from inventory.json

for item_id, predicted_qty in predictions.items():
    recipe = recipes.get(item_id, {})
    for ing in recipe.get("ingredients", []):
        needed = ing["qty"] * predicted_qty
        current = get_stock(ing["ingredient_id"])
        if current < needed:
            shortfall = needed - current
            alert(f"Order {shortfall} {ing['unit']} of {ing['name']}")
\end{lstlisting}


\subsection{Integration with Petpooja}
\begin{itemize}[nosep]
    \item Create API endpoint: \texttt{GET /forecast/\{item\_id\}?days=7}
    \item Run daily cron job to retrain models at 2:00 AM
    \item Display predictions on the Dashboard with confidence bands (chart)
    \item Auto-generate purchase order suggestions when predicted demand exceeds stock
\end{itemize}


\subsection{Evaluation Metrics}
\begin{center}
\begin{tabular}{lll}
\toprule
\textbf{Metric} & \textbf{Description} & \textbf{Target} \\
\midrule
MAE   & Mean Absolute Error per item/day & $< 3$ units \\
MAPE  & Mean Absolute Percentage Error   & $< 15\%$ \\
Coverage & \% of actuals within prediction interval & $> 90\%$ \\
\bottomrule
\end{tabular}
\end{center}


\newpage
% ════════════════════════════════════════════════════════════════════
\section{NLP Intent Parser (LLM-Based Order Understanding)}
% ════════════════════════════════════════════════════════════════════

\begin{infobox}[Model Overview]
\begin{tabularx}{\textwidth}{@{}lX@{}}
\textbf{Goal}       & Replace regex + fuzzy-match parsing with an LLM that understands complex natural-language orders \\
\textbf{Algorithm}  & Large Language Model (prompt engineering / fine-tuning) \\
\textbf{Library}    & \texttt{transformers} + \texttt{torch} (local) \textbf{or} \texttt{openai} SDK (cloud) \\
\textbf{Priority}   & \textcolor{red}{\textbf{HIGH}} --- dramatically improves voice order accuracy \\
\end{tabularx}
\end{infobox}


\subsection{Why LLM Over Regex + Fuzzy Match?}

The current system uses RapidFuzz for fuzzy string matching, which works well for simple orders but fails on:

\begin{itemize}[nosep]
    \item \textbf{Corrections:} ``Give me 2 burgers, actually make it 3''
    \item \textbf{References:} ``Same thing as last order but no onions''
    \item \textbf{Multi-item modifiers:} ``One margherita and one pepperoni, extra cheese on both''
    \item \textbf{Quantities in words:} ``A couple of cold coffees and half a dozen samosas''
    \item \textbf{Hinglish/slang:} ``Ek paneer wala pizza aur do chai de do''
\end{itemize}

An LLM understands context, intent, and corrections naturally.


\subsection{Implementation Strategy}

\subsubsection{Option A: Local Model (Google Gemma 2B)}

This runs entirely on your machine --- no API cost, no data leaving your server.

\begin{lstlisting}[language=bash]
pip install transformers torch accelerate
\end{lstlisting}

\begin{lstlisting}
from transformers import AutoTokenizer, AutoModelForCausalLM
import json

# Load model (first run downloads ~5 GB)
model_name = "google/gemma-2b-it"
tokenizer = AutoTokenizer.from_pretrained(model_name)
model = AutoModelForCausalLM.from_pretrained(model_name)

def parse_order_llm(transcription: str, menu_items: list) -> dict:
    menu_str = "\n".join(
        f"- {m['item_id']}: {m['name']} (Rs {m['selling_price']})"
        for m in menu_items
    )
    
    prompt = f"""You are a restaurant order parser.
Given the spoken order and menu, extract items as JSON.

MENU:
{menu_str}

CUSTOMER ORDER: "{transcription}"

Respond ONLY with valid JSON:
{{"items": [{{"item_id": "...", "name": "...", "qty": N,
              "modifiers": [...], "price": N}}], 
  "total": N}}"""
    
    inputs = tokenizer(prompt, return_tensors="pt")
    outputs = model.generate(**inputs, max_new_tokens=512,
                             temperature=0.1)
    response = tokenizer.decode(outputs[0], skip_special_tokens=True)
    
    # Extract JSON from response
    json_start = response.rfind("{")
    json_end = response.rfind("}") + 1
    return json.loads(response[json_start:json_end])
\end{lstlisting}


\subsubsection{Option B: Cloud API (OpenAI GPT-4o-mini)}

Lower latency ($<1$s), higher accuracy, but costs \$0.15 per 1M input tokens.

\begin{lstlisting}
import openai, json

client = openai.OpenAI(api_key="sk-...")

def parse_order_api(transcription: str, menu_items: list) -> dict:
    menu_str = "\n".join(
        f"- {m['item_id']}: {m['name']} (Rs {m['selling_price']})"
        for m in menu_items
    )
    
    response = client.chat.completions.create(
        model="gpt-4o-mini",
        response_format={"type": "json_object"},
        messages=[
            {"role": "system",
             "content": f"You parse restaurant orders.\nMENU:\n{menu_str}"},
            {"role": "user", "content": transcription}
        ],
        temperature=0.1
    )
    return json.loads(response.choices[0].message.content)
\end{lstlisting}


\subsection{Fine-Tuning (Optional --- Advanced)}

For the highest accuracy, fine-tune a small model on your own data.

\subsubsection{Step 1: Collect Training Data}
Create a JSONL file with 100--500 labeled examples:
\begin{lstlisting}
{"input": "2 paneer burgers extra cheese and a cold coffee",
 "output": {"items": [
   {"item_id": "M002", "name": "Paneer Burger", "qty": 2,
    "modifiers": ["extra cheese"]},
   {"item_id": "M027", "name": "Cold Coffee", "qty": 1,
    "modifiers": []}
 ]}}
\end{lstlisting}

\subsubsection{Step 2: Fine-Tune with QLoRA}
\begin{lstlisting}[language=bash]
pip install peft trl datasets bitsandbytes
\end{lstlisting}

QLoRA allows fine-tuning a 2B parameter model on a single GPU (8~GB VRAM) by quantizing the base model to 4-bit and training only small adapter layers.

\begin{itemize}[nosep]
    \item Training data: 100--500 labeled examples
    \item Training time: $\sim$30 minutes on a single GPU
    \item Expected result: 95\%+ accuracy on order parsing
\end{itemize}


\subsection{Evaluation Metrics}
\begin{center}
\begin{tabular}{lll}
\toprule
\textbf{Metric} & \textbf{Description} & \textbf{Target} \\
\midrule
Exact Match Rate  & \% of orders with 100\% correct items & $> 90\%$ \\
Item-Level F1     & Precision/recall for item extraction   & $> 0.95$ \\
Modifier Accuracy & \% of modifiers correctly captured     & $> 85\%$ \\
Latency           & End-to-end response time               & $< 2$s (local), $< 1$s (API) \\
\bottomrule
\end{tabular}
\end{center}


\newpage
% ════════════════════════════════════════════════════════════════════
\section{Collaborative Filtering (Personalized Upsell)}
% ════════════════════════════════════════════════════════════════════

\begin{infobox}[Model Overview]
\begin{tabularx}{\textwidth}{@{}lX@{}}
\textbf{Goal}       & Learn customer ordering patterns to suggest personalized upsells \\
\textbf{Algorithm}  & SVD (Matrix Factorization) / Neural Collaborative Filtering \\
\textbf{Library}    & \texttt{scikit-surprise} (SVD) or \texttt{tensorflow} (Neural CF) \\
\textbf{Priority}   & \textcolor{orange}{\textbf{MEDIUM}} --- requires customer identity tracking \\
\end{tabularx}
\end{infobox}


\subsection{How It Works}

Collaborative Filtering (CF) builds a \textbf{user--item interaction matrix} from order history. It discovers that ``customers who ordered Paneer Pizza also loved Garlic Bread'' and generates personalized recommendations.

\begin{enumerate}[nosep]
    \item Build matrix: rows = customers, columns = menu items, values = order frequency
    \item Factorize the matrix into latent factors using SVD
    \item For a given customer, predict scores for un-ordered items
    \item Return top-K items with highest predicted scores
\end{enumerate}

\subsubsection{User--Item Matrix Example}
\begin{center}
\begin{tabular}{lccccc}
\toprule
              & M001 & M003 & M009 & M019 & M027 \\
\midrule
Customer\_001 & 5    & 0    & 3    & 4    & 2    \\
Customer\_002 & 0    & 7    & 1    & 6    & 0    \\
Customer\_003 & 3    & 2    & 0    & 1    & 5    \\
\bottomrule
\end{tabular}
\end{center}
\textit{Values = number of times customer ordered that item.}


\subsection{Data Requirements}
\begin{itemize}[nosep]
    \item Customer identification: phone number, loyalty ID, or table session
    \item Minimum: \textbf{500 unique customers} with \textbf{3+ orders each}
    \item More data = better recommendations (cold-start is the main challenge)
\end{itemize}


\subsection{Step-by-Step Training}

\subsubsection{Step 1: Install Dependencies}
\begin{lstlisting}[language=bash]
pip install scikit-surprise pandas numpy
\end{lstlisting}

\subsubsection{Step 2: Prepare Interaction Data}
\begin{lstlisting}
import pandas as pd
from surprise import Dataset, Reader, SVD
from surprise.model_selection import cross_validate

# Build user-item-rating dataframe from orders
rows = []
for order in orders:
    cust_id = order.get("customer_id", "anonymous")
    for item in order["items"]:
        rows.append({
            "user": cust_id,
            "item": item["item_id"],
            "rating": item["qty"]   # frequency as implicit rating
        })

df = pd.DataFrame(rows)
df = df.groupby(["user", "item"])["rating"].sum().reset_index()
df["rating"] = df["rating"].clip(1, 5)  # normalize to 1-5 scale
\end{lstlisting}

\subsubsection{Step 3: Train SVD Model}
\begin{lstlisting}
reader = Reader(rating_scale=(1, 5))
data = Dataset.load_from_df(df[["user", "item", "rating"]], reader)

# Train SVD
algo = SVD(n_factors=50, n_epochs=30, lr_all=0.005, reg_all=0.02)
trainset = data.build_full_trainset()
algo.fit(trainset)

# Cross-validate
results = cross_validate(algo, data, measures=["RMSE", "MAE"], cv=5)
print(f"RMSE: {results['test_rmse'].mean():.3f}")
print(f"MAE:  {results['test_mae'].mean():.3f}")
\end{lstlisting}

\subsubsection{Step 4: Generate Recommendations}
\begin{lstlisting}
def get_recommendations(customer_id, menu_items, top_k=5):
    """Get top-K item recommendations for a customer."""
    ordered = set(df[df["user"] == customer_id]["item"])
    candidates = [m["item_id"] for m in menu_items
                  if m["item_id"] not in ordered]
    
    predictions = [
        (iid, algo.predict(customer_id, iid).est)
        for iid in candidates
    ]
    predictions.sort(key=lambda x: x[1], reverse=True)
    return predictions[:top_k]
\end{lstlisting}

\subsubsection{Step 5: Cold-Start Fallback}
\begin{lstlisting}
def get_popular_items(top_k=5):
    """Fallback for new customers with no order history."""
    popularity = df.groupby("item")["rating"].sum()
    return popularity.nlargest(top_k).index.tolist()
\end{lstlisting}


\subsection{Integration with Petpooja}
\begin{itemize}[nosep]
    \item Create endpoint: \texttt{GET /recommend/\{customer\_id\}?top\_k=5}
    \item Add \texttt{customer\_id} field to the order flow (phone number or loyalty card)
    \item Replace generic Apriori upsells with personalized CF recommendations
    \item Fall back to popularity-based suggestions for new customers
    \item Weekly retraining via cron job with latest order data
\end{itemize}


\subsection{Evaluation Metrics}
\begin{center}
\begin{tabular}{lll}
\toprule
\textbf{Metric} & \textbf{Description} & \textbf{Target} \\
\midrule
RMSE       & Root Mean Square Error on held-out data     & $< 1.0$ \\
Hit Rate@5 & \% of test orders in top-5 recommendations & $> 30\%$ \\
NDCG@5     & Normalized Discounted Cumulative Gain      & $> 0.4$ \\
Coverage   & \% of catalog items ever recommended       & $> 60\%$ \\
\bottomrule
\end{tabular}
\end{center}


\newpage
% ════════════════════════════════════════════════════════════════════
\section{Anomaly Detection (Revenue \& Inventory)}
% ════════════════════════════════════════════════════════════════════

\begin{infobox}[Model Overview]
\begin{tabularx}{\textwidth}{@{}lX@{}}
\textbf{Goal}       & Detect unusual patterns in revenue, orders, and inventory in real-time \\
\textbf{Algorithm}  & Z-Score (statistical) + Isolation Forest (ML) \\
\textbf{Library}    & \texttt{scikit-learn}, \texttt{numpy}, \texttt{pandas} \\
\textbf{Priority}   & \textcolor{orange}{\textbf{MEDIUM}} --- critical for loss prevention \\
\end{tabularx}
\end{infobox}


\subsection{Types of Anomalies Detected}
\begin{itemize}[nosep]
    \item \textbf{Revenue:} Sudden revenue drop $> 30\%$ compared to same day last week
    \item \textbf{Orders:} Unusually high void/cancel rate (potential fraud)
    \item \textbf{Inventory:} Ingredients depleting faster than order volume justifies (theft/waste)
    \item \textbf{Timing:} Orders placed outside business hours (system misuse)
    \item \textbf{Pricing:} Unauthorized discount patterns
\end{itemize}


\subsection{Data Requirements}
\begin{itemize}[nosep]
    \item Hourly revenue totals --- at least 14 days (30+ recommended)
    \item Order-level data with timestamps, void status, discount flags
    \item Inventory movement logs (depletion events with timestamps)
\end{itemize}


\subsection{Step-by-Step Training}

\subsubsection{Step 1: Install Dependencies}
\begin{lstlisting}[language=bash]
pip install scikit-learn pandas numpy
\end{lstlisting}

\subsubsection{Step 2: Statistical Anomaly Detection (Z-Score)}

Simple but effective for single-metric monitoring (e.g., daily revenue):

\begin{lstlisting}
import numpy as np
import pandas as pd

def detect_revenue_anomalies(daily_revenue, threshold=2.5):
    """Flag days where revenue deviates > 2.5 std from mean."""
    mean = daily_revenue.mean()
    std = daily_revenue.std()
    z_scores = (daily_revenue - mean) / std
    
    anomalies = daily_revenue[abs(z_scores) > threshold]
    return [
        {
            "date": str(date),
            "revenue": float(val),
            "z_score": float(z),
            "type": "low" if z < 0 else "high"
        }
        for date, val, z in zip(
            anomalies.index, anomalies.values,
            z_scores[abs(z_scores) > threshold]
        )
    ]
\end{lstlisting}


\subsubsection{Step 3: Isolation Forest (Multi-Dimensional)}

For detecting complex anomalies across multiple features simultaneously:

\begin{lstlisting}
from sklearn.ensemble import IsolationForest
from sklearn.preprocessing import StandardScaler

# Feature matrix: each row = 1 hour of operations
# Columns: [revenue, order_count, avg_order_value,
#            void_rate, discount_rate, inventory_depletion]
features = build_feature_matrix(orders, inventory_logs)

# Scale features
scaler = StandardScaler()
X_scaled = scaler.fit_transform(features)

# Train Isolation Forest
model = IsolationForest(
    n_estimators=200,
    contamination=0.05,   # expect ~5% anomalies
    max_features=0.8,
    random_state=42
)
model.fit(X_scaled)

# Predict: -1 = anomaly, 1 = normal
predictions = model.predict(X_scaled)
scores = model.decision_function(X_scaled)
print(f"Anomalies found: {(predictions == -1).sum()}")
\end{lstlisting}


\subsubsection{Step 4: Real-Time Scoring}
\begin{lstlisting}
import pickle

# Save model + scaler for production use
with open("models/anomaly_model.pkl", "wb") as f:
    pickle.dump({"model": model, "scaler": scaler}, f)

def score_current_hour(current_features):
    """Score the current hour's data for anomalies."""
    with open("models/anomaly_model.pkl", "rb") as f:
        bundle = pickle.load(f)
    
    X = bundle["scaler"].transform([current_features])
    is_anomaly = bundle["model"].predict(X)[0] == -1
    score = bundle["model"].decision_function(X)[0]
    
    if is_anomaly:
        return {
            "alert": True,
            "severity": "HIGH" if score < -0.3 else "MEDIUM",
            "score": float(score),
            "message": classify_anomaly(current_features)
        }
    return {"alert": False}
\end{lstlisting}


\subsubsection{Step 5: Root Cause Classification}
\begin{lstlisting}
def classify_anomaly(features):
    """Determine the most likely root cause."""
    revenue, orders, aov, voids, discounts, depletion = features
    
    if voids > historical_void_mean * 2:
        return "High void rate - possible fraud"
    if depletion > expected_depletion * 1.5:
        return "Excessive inventory depletion - check for waste/theft"
    if revenue < historical_revenue_mean * 0.6:
        return "Significant revenue drop - investigate operations"
    if discounts > historical_discount_mean * 2:
        return "Unusual discount pattern - verify authorization"
    return "General anomaly - manual review needed"
\end{lstlisting}


\subsection{Integration with Petpooja}
\begin{itemize}[nosep]
    \item Create endpoint: \texttt{GET /alerts/anomalies?hours=24}
    \item Run scoring every hour via APScheduler background task
    \item Display alerts on Dashboard with severity badges (HIGH / MEDIUM / LOW)
    \item Send push notifications for HIGH severity anomalies
    \item Weekly model retraining on a rolling 60-day window
\end{itemize}


\subsection{Evaluation Metrics}
\begin{center}
\begin{tabular}{lll}
\toprule
\textbf{Metric} & \textbf{Description} & \textbf{Target} \\
\midrule
Precision        & \% of flagged anomalies that are real     & $> 80\%$ \\
Recall           & \% of actual anomalies caught             & $> 90\%$ \\
False Positive Rate & Avoid alert fatigue                   & $< 5\%$ \\
Time to Detect   & Anomaly flagged within N hours           & $< 1$ hour \\
\bottomrule
\end{tabular}
\end{center}


\newpage
% ════════════════════════════════════════════════════════════════════
\section{Summary \& Technology Stack}
% ════════════════════════════════════════════════════════════════════

\subsection{Model Comparison}

\begin{center}
\begin{tabular}{lllll}
\toprule
\textbf{Model} & \textbf{Algorithm} & \textbf{Library} & \textbf{Min Data} & \textbf{Priority} \\
\midrule
Demand Forecasting   & Prophet / LSTM       & \texttt{prophet}    & 30 days orders     & HIGH \\
NLP Intent Parser    & LLM (Gemma/GPT)      & \texttt{transformers} & Menu + prompt    & HIGH \\
Collaborative Filter & SVD / Neural CF      & \texttt{surprise}   & 500 customers      & MEDIUM \\
Anomaly Detection    & Isolation Forest     & \texttt{scikit-learn} & 14 days ops data & MEDIUM \\
\bottomrule
\end{tabular}
\end{center}


\subsection{Full Technology Stack}

\begin{itemize}[nosep]
    \item \textbf{Core:} Python 3.12, FastAPI, pandas, numpy
    \item \textbf{Speech:} OpenAI Whisper (already implemented)
    \item \textbf{NLP:} \texttt{transformers} + \texttt{torch} (local) or \texttt{openai} SDK (cloud)
    \item \textbf{Forecasting:} \texttt{prophet} (Facebook) or \texttt{tensorflow}
    \item \textbf{Recommendations:} \texttt{scikit-surprise}
    \item \textbf{Anomaly Detection:} \texttt{scikit-learn} (Isolation Forest)
    \item \textbf{Scheduling:} APScheduler or Celery for periodic retraining
    \item \textbf{Model Storage:} Pickle files or MLflow for versioning
    \item \textbf{Hardware:} GPU recommended for NLP model (NVIDIA T4 minimum)
\end{itemize}


\subsection{Implementation Roadmap}

\begin{enumerate}
    \item \textbf{Phase 1 (Weeks 1--2):} Implement Demand Forecasting with Prophet. Connect predictions to inventory alerts on the Dashboard.
    \item \textbf{Phase 2 (Weeks 3--4):} Integrate NLP Intent Parser. Start with cloud API (GPT-4o-mini) for quick results, migrate to local Gemma model later.
    \item \textbf{Phase 3 (Weeks 5--8):} Deploy Collaborative Filtering once customer identity tracking is in place. Requires loyalty or phone-number based identification.
    \item \textbf{Phase 4 (Weeks 9--12):} Add Anomaly Detection. Start with Z-Score for revenue monitoring, then add Isolation Forest for multi-dimensional analysis.
\end{enumerate}

\vfill
\begin{center}
{\color{orange}\rule{8cm}{1.5pt}}\\[8pt]
{\color{gray}\textit{Petpooja AI Copilot --- ML Model Implementation Guide}}\\[4pt]
{\color{gray}\textit{Auto-generated $\cdot$ \today}}
\end{center}

\end{document}
